% ~10 min
\partframe{Work in Progress}{Down to Low-Level Code}

\begin{frame}
  \frametitle{Two Ways Out of Rocq}
  \begin{center}
    \begin{tikzpicture}[alt={A specification is calculated by laws into
        a high-level program. Today, the high-level program is extracted to
        OCaml, relying on garbage-collected lists and trees, closures and
        boxed integers. The goal: more laws take the high-level program down
        to a low-level form, which is printed directly as low-level code
        with flat buffers, loops, machine integers, no allocation, and
        parallel hardware. The specification, the high-level program and the
        low-level form sit in a box labelled verified in Rocq.},
      stage/.style = {box, minimum height = 10mm, text width = 24mm,
        align = center},
      lbl/.style = {font = \bfseries},
      rocq/.style = {draw = accent, fill = accent!6, rounded corners = 6pt,
        thick, inner xsep = 3mm, inner ysep = 6mm}]
      \node[stage] (s) at (0, 0) {specification};
      \node[stage, draw = good, fill = good!8] (h) at (4.6, 0)
        {high-level program};
      \node[stage, right = 20mm of h] (o) {OCaml};
      \node[lbl, above = 1mm of o] {today};
      \draw[->, thick] (s) -- node[above, note] {laws} (h);
      \draw[->, thick] (h) -- node[pos = .6, above, note] {extract} (o);
      \node[note, right = 3mm of o, align = left, text width = 28mm]
        {GC'd lists, trees,\\ closures,\\ boxed integers};
      \uncover<2->{
        \node[stage, draw = good, fill = good!8] (lf) at (4.6, -2.6)
          {low-level form};
        \node[hot, minimum height = 10mm, text width = 24mm,
          align = center, right = 20mm of lf] (l) {\textbf{low-level code}};
        \node[lbl, above = 1mm of l] {goal};
        \draw[->, thick] (h) -- node[right, note] {laws} (lf);
        \draw[->, thick] (lf) -- node[pos = .6, above, note] {print} (l);
        \node[note, right = 3mm of l, align = left, text width = 28mm]
          {flat buffers, loops,\\ machine integers,\\
          no allocation,\\ parallel hardware};
      }
      \begin{scope}[on background layer]
        \only<1>{\node[rocq, fit = (s) (h)] (r) {};}
        \only<2->{\node[rocq, fit = (s) (h) (lf)] (r) {};}
        \node[note, text = accent, anchor = south west, font = \small\bfseries]
          at (r.south west) {\ verified in Rocq};
      \end{scope}
    \end{tikzpicture}
  \end{center}
  \uncover<3->{Not ``OCaml, then lower'': \textbf{more calculation} reaches the
    low-level form; extraction is a direct print.}
  \speakernote{What we have today exits Rocq through OCaml extraction:
    idiomatic and safe, but it relies on a garbage collector, linked data
    structures, closures, boxed integers. [click] Fast code for numerical
    kernels, cryptography, embedded systems, ML runtimes looks different:
    data in flat buffers, loops instead of recursion, machine integers, no
    allocation in the hot path, work split across cores or vector units.
    The aim is to get there directly: calculate, inside Rocq, into a
    representation that is already low-level, and print that, without
    going through OCaml. [click] So the low-level decisions (where data
    lives, how work is scheduled) are calculation steps, proved like the
    others, not a separate compiler pass. Today this step is done by hand
    (and is where bugs live), or not at all: a compiler will not change
    your data layout or parallelise your algorithm for you.}
\end{frame}

\begin{frame}
  \frametitle{What We Are Working On}
  \vfill
  \begin{center}
    \begin{minipage}{.82\textwidth}
      \begin{sticky}[halign = center]
        {\Large \textbf{Low-level}, \textbf{correct-by-construction} code,
        extracted from \textbf{high-level calculation steps},\par}
        \medskip
        {\large possibly guided by a \textbf{cost model}.\par}
      \end{sticky}
    \end{minipage}
  \end{center}
  \vfill
  \begin{itemize}[<+->]
    \item You choose \emph{laws}, not pointers and indices.
    \item Every step proved; the output is machine-level code.
    \item A cost model says \emph{which} equal program to pick.
  \end{itemize}
  \speakernote{This is the work in progress: the direction, not
    finished results. The programmer, or a tool, works
    at the level of the laws we have seen: fuse this, split that, memoise
    here. The output is not OCaml but low-level code: explicit memory,
    loops, machine arithmetic. Every step between the two is a proved
    equation, so the low-level code is correct by construction, and no
    hand-written proof about pointers or indices is needed. Many programs
    are equal to the specification; which one is best depends on the
    machine, so the choice should be guided by a cost model. Next: three
    pieces we are working on.}
\end{frame}

\begin{frame}
  \frametitle{Ongoing: Laws That Change Representation}
  \begin{center}
    \begin{tikzpicture}[alt={Three kinds of rewrite. A tree of
        intermediate results becomes a flat buffer. Recursion becomes a
        loop. One traversal becomes independent pieces that can run at the
        same time.},
      from/.style = {box, text width = 32mm, align = center,
        minimum height = 10mm},
      to/.style = {box, draw = good, fill = good!8, text width = 32mm,
        align = center, minimum height = 10mm}]
      \foreach \y/\f/\t/\o in {
          0/{intermediate tree}/{flat buffer}/1,
          -1.6/{recursion}/{loop}/2,
          -3.2/{one traversal}/{independent pieces, in parallel}/3} {
        \uncover<\o->{
          \node[from] (f\o) at (0, \y) {\f};
          \node[to, right = 30mm of f\o] (t\o) {\t};
          \draw[->, thick] (f\o) -- node[above, note] {$=$, proved once}
            (t\o);
        }
      }
    \end{tikzpicture}
  \end{center}
  \uncover<4->{Today's laws change the \emph{algorithm}; these change
    \emph{where data lives} and \emph{how work is scheduled}.}
  \speakernote{The laws we have remove or add intermediate structures, but
    the only intermediate structure they know is a tree. Low-level code
    needs laws whose right-hand side is a flat buffer in memory. [click]
    Divide and conquer with a known shape can become a loop. [click] When
    pieces of the work are independent (for instance, they write to
    disjoint parts of a buffer), they can run at the same time; that
    independence should itself be a proved property of the split, once,
    not per program. [click] Same style as before: general equations,
    proved once, used as rewrite steps.}
\end{frame}

\begin{frame}
  \frametitle{Ongoing: A Target the Laws Never Leave}
  \begin{center}
    \begin{tikzpicture}[alt={Nested regions: all programs, containing a
        low-level fragment. Calculation steps move the specification, which
        starts outside, into the low-level fragment and keep it there. A
        small trusted printer turns programs in the fragment into C or
        assembly.},
      region/.style = {draw, rounded corners = 6pt, thick}]
      \node[region, draw = blue!40, fill = blue!4, minimum width = 100mm,
        minimum height = 44mm] (all) at (0, 0) {};
      \node[note, anchor = north west] at (all.north west) {\ all programs};
      \node[region, draw = good, fill = good!8, minimum width = 48mm,
        minimum height = 30mm] (low) at (22mm, -3mm) {};
      \node[note, anchor = north west, text = good] at (low.north west)
        {\ low-level fragment};
      \node[box] (s) at (-36mm, 6mm) {spec};
      \node[box] (m) at (-12mm, -6mm) {$\cdot$};
      \node[box] (p1) at (12mm, 0mm) {$\cdot$};
      \node[box] (p2) at (32mm, -10mm) {$\cdot$};
      \draw[->, thick] (s) -- (m);
      \draw[->, thick] (m) -- (p1);
      \draw[->, thick] (p1) -- (p2);
      \uncover<2->{
        \node[hot, right = 18mm of all.east, anchor = west] (c)
          {C / assembly};
        \draw[->, thick] (all.east) ++(0, -3mm) -- node[above, note]
          {print} (c.west);
      }
    \end{tikzpicture}
  \end{center}
  \begin{itemize}
    \item Every law keeps programs \textbf{inside} the fragment.
    \item<2-> Only a \textbf{small printer} is trusted.
    \item<3-> Overflow and bounds: \textbf{side conditions} of the laws.
  \end{itemize}
  \speakernote{Extraction to OCaml is a large trusted step. For low-level
    code we want a first-order fragment (no closures, explicit memory,
    machine integers) such that the laws, applied to programs in the
    fragment, stay in it. Then the calculation moves the specification
    into the fragment, and keeps it there. [click] Printing a program in
    that fragment as C or assembly is almost syntactic, so the trusted
    base is small. [click] The laws are true for mathematical integers and
    unbounded memory; the code runs with 64-bit integers and fixed buffers.
    Those facts become proof obligations attached to the steps, ideally
    discharged automatically so they do not get in the way of rewriting.}
\end{frame}

\begin{frame}
  \frametitle{Ongoing: A Cost Model}
  \begin{center}
    \begin{tikzpicture}[alt={One specification with two equal programs
        below it. The fused program uses no intermediate memory but runs
        sequentially. The split program writes an intermediate buffer but
        its two halves can run in parallel or be vectorised. Which is
        faster depends on the machine and the input size.},
      prog/.style = {box, draw = good, fill = good!8, text width = 40mm,
        align = center, minimum height = 14mm}]
      \node[box, minimum height = 9mm] (s) at (0, 0) {specification};
      \node[prog] (f) at (-34mm, -26mm) {\textbf{fused}\\ \small no buffer,
        one core};
      \node[prog] (p) at (34mm, -26mm) {\textbf{split}\\ \small buffer,
        many cores / SIMD};
      \draw[->, thick] (s) -- node[left, note] {$=$} (f);
      \draw[->, thick] (s) -- node[right, note] {$=$} (p);
      \uncover<2->{
        \node[hot, below = 6mm of f] {cost: memory, latency};
        \node[hot, below = 6mm of p] {cost: bandwidth, sync};
      }
    \end{tikzpicture}
  \end{center}
  \uncover<3->{The laws say what is \emph{equal}; a cost model says what is
    \emph{better}.}
  \speakernote{Deforestation can be used in both directions: fuse two
    passes, or split one pass through an intermediate buffer. Both are
    equal to the specification. [click] Fused saves memory traffic; split
    may enable parallelism or vectorisation, at the price of a buffer and
    synchronisation. Which is better depends on the hardware and the input
    size. [click] The laws cannot see that. A cost model over the low-level
    fragment can, and can rank the candidate derivations.}
\end{frame}

\begin{frame}
  \frametitle{Untrusted Search, Trusted Checking}
  \begin{center}
    \begin{tikzpicture}[alt={A specification and a cost model go into a
        search, which may use heuristics, equality saturation or a learned
        model, and is not trusted. It proposes a chain of laws. Rocq checks
        the chain, and the result is printed as low-level code. A wrong
        proposal is rejected.},
      stage/.style = {minimum height = 12mm, text width = 24mm,
        align = center}]
      \node[box, stage] (s) {specification\\ + cost model};
      \node[box, stage, right = 8mm of s, draw = bad, fill = bad!6] (q)
        {\textbf{search}\\ \small heuristics, e-graphs, ML};
      \node[box, stage, right = 22mm of q, draw = good, fill = good!8] (r)
        {\textbf{Rocq checks}\\ \small each step};
      \node[hot, stage, right = 8mm of r] (o) {low-level code};
      \draw[->, thick] (s) -- (q);
      \draw[->, thick] (q) -- node[above, note, align = center]
        {chain of\\ laws} (r);
      \draw[->, thick] (r) -- (o);
      \node[note, below = 2mm of q, text = bad] {untrusted};
      \node[note, below = 2mm of r, text = good] {trusted, small};
      \uncover<2->{
        \draw[->, thick, bad, dashed] (r.north) to[bend right = 30]
          node[above, note, text = bad] {rejected: try again} (q.north);
      }
    \end{tikzpicture}
  \end{center}
  \begin{itemize}
    \item The derivation is a \textbf{certificate}.
    \item<2-> Search can be as clever, or as wrong, as it likes.
  \end{itemize}
  \speakernote{How the pieces fit. Finding a good derivation is a search
    problem over equal programs, guided by the cost model: hand-written
    heuristics, equality saturation with e-graphs, or a learned model
    proposing the next rewrite. None of that needs to be trusted: what it
    produces is a chain of laws, which Rocq checks. [click] If a proposed
    step is wrong, the check fails and the search tries again; the output
    is correct by construction no matter how it was found. For ML people:
    the same pattern as a model proposing proofs to a proof checker. For
    security people: the trusted base is the kernel and a small printer,
    not the optimiser.}
\end{frame}

\begin{frame}
  \frametitle{Where We Are}
  \begin{center}
    \begin{tikzpicture}[alt={An index range from 0 to n is split into four
        pieces of different lengths. The pieces cover the range exactly,
        with no gaps and no overlaps; this is proved.},
      piece/.style = {draw = good, fill = good!8, minimum height = 7mm,
        inner sep = 0pt}]
      \node[piece, minimum width = 30mm, anchor = west] (p0) at (0, 0) {};
      \node[piece, minimum width = 18mm, right = 0mm of p0] (p1) {};
      \node[piece, minimum width = 24mm, right = 0mm of p1] (p2) {};
      \node[piece, minimum width = 26mm, right = 0mm of p2] (p3) {};
      \node[note, below = 1mm of p0.south west] {$0$};
      \node[note, below = 1mm of p3.south east] {$n$};
      \node[note, right = 3mm of p3] {exactly covered, no overlaps};
    \end{tikzpicture}
  \end{center}
  \begin{columns}[t]
    \begin{column}{0.5\textwidth}
      \textbf{\textcolor{good}{In Rocq, axiom-free:}}
      \begin{itemize}
        \item Splitting index ranges $k$ ways
        \item Machine integers, no overflow (proved)
        \item Buffers, compared by their reads
        \item Tree reductions over ranges
        \item Writes to a buffer, driven by any split
      \end{itemize}
    \end{column}
    \begin{column}{0.45\textwidth}
      \textbf{\textcolor{gray}{Not yet:}}
      \begin{itemize}
        \item Laws that add or remove a buffer
        \item Each index written exactly once
        \item The low-level target and printer
        \item The cost model and search
      \end{itemize}
    \end{column}
  \end{columns}
  \speakernote{Honest status: early. In Rocq, without axioms: index ranges
    split into k pieces, with a proof that the pieces cover the range
    exactly with no overlap (the basis for exactly-once writes and for
    independent pieces). Indices are abstract, so the same code works for
    unbounded naturals and for machine integers, with overflow as explicit
    side conditions; extraction specialises to native integers. Buffers
    with an equality that compares reads, so no axioms. A tree reduction
    over a range as a hylomorphism. Writes into a buffer as a monoid
    algebra driven by any terminating split, so the fusion law applies.
    Not yet: the laws that introduce or remove a buffer, the exactly-once
    property, the low-level target, the cost model and search.}
\end{frame}

\begin{frame}
  \frametitle{Summary}
  \begin{itemize}
    \item Program optimisation as \textbf{algebra}: rewrite with laws.
    \item Divide and conquer as \textbf{hylomorphisms}: laws proved once.
    \item In Rocq: \textbf{machine-checked}, no extra axioms, clean OCaml.
    \item Ongoing: \textbf{low-level, correct-by-construction code} from
      high-level steps, guided by a \textbf{cost model}.
  \end{itemize}
  \speakernote{Recap. The calculation style is old; the contribution is
    making it work in a proof assistant without extra axioms and with
    runnable output. The ongoing work is to push the output down to low-level
    code, with the cost model choosing between equal programs and an
    untrusted search proposing the steps.}
\end{frame}

\begin{frame}
  \vfill
  \begin{center}
    \begin{minipage}{.6\textwidth}
      \begin{sticky}[halign = center]
        {\Huge Thank you!\par}
        \bigskip
        {\large Questions?\par}
        \medskip
        \url{https://github.com/dcastrop/rocq-hylomorphisms}
        \medskip

        {\small Based on joint work with Marco Paviotti and Michael
        Vollmer.\par}
      \end{sticky}
    \end{minipage}
  \end{center}
  \vfill
  \speakernote{Based on joint work with Marco and Michael. Questions.}
\end{frame}
